\documentclass[a4paper,14pt]{extarticle}
\usepackage[T2A]{fontenc}
\usepackage[utf8]{inputenc}
\usepackage[russian]{babel}
\usepackage{tempora,setspace,indentfirst,tabularx,titlesec}
\usepackage[left=30mm,right=15mm,top=20mm,bottom=20mm]{geometry}
\onehalfspacing
\setlength{\parindent}{1.25cm}
\titlespacing*{\section}{\parindent}{12pt}{6pt}
\begin{document}
\section*{Проект АРМ обработки заявок}
Учебный пример. Проектируется прототип для одного отдела.
Оператор регистрирует и обрабатывает заявки; наблюдатель только читает их.
\par\vspace{6pt}

\noindent
\begin{tabularx}{\linewidth}{|l|>{\raggedright\arraybackslash}X|>{\raggedright\arraybackslash}X|}
\hline
ID & Требование & Проверка \\
\hline
F1 & Регистрация заявки & Пустая тема отклоняется; корректная заявка получает номер и автора. \\
\hline
F2 & Новая $\to$ В работе $\to$ Закрыта & Для начала нужен исполнитель, для закрытия --- результат. \\
\hline
F3 & Поиск и фильтр & Карточка найдена по номеру; фильтр оставляет нужное состояние. \\
\hline
F4 & Наблюдатель только читает & Сервер отклоняет изменение; данные не меняются. \\
\hline
F5 & Согласованное сохранение & История сохраняется с заявкой; устаревшая версия отклоняется. \\
\hline
\end{tabularx}

\section*{Модель данных}
Пользователь: номер, логин, роль, признак активности.
Заявка: номер, тема, описание, состояние, время создания, автор,
исполнитель, результат и версия. История: ссылка на заявку,
старое и новое состояния, автор и время изменения.

При регистрации состояние --- «Новая», старое состояние в истории пусто.
Изменение состояния и история сохраняются в одной транзакции.
Устаревшая версия карточки требует обновления данных перед повторной попыткой.

\section*{Границы проверки}
Для каждого требования оформляется протокол с ожидаемым и фактическим результатом.
Нагрузку и восстановление после сбоя проверяют отдельно.
\end{document}
